\documentclass[10pt,a4paper]{amsart}
\usepackage[margin=19mm]{geometry}
\usepackage{amsmath,amssymb,mathtools}
\usepackage[scr=rsfs,cal=euler]{mathalfa}
\usepackage{xfrac}
\usepackage[unicode,hidelinks,bookmarks=false]{hyperref}
\newcommand{\ie}{i.e.\ }
\newcommand{\cf}{cf.\ }
\newcommand{\resp}{respectively\ }
\newcommand{\Subsection}{\subsection}
\providecommand{\nobreakdash}{\nobreak\hbox{-}}
\setlength{\emergencystretch}{2em}
\multlinegap0pt
\allowdisplaybreaks[4]
\newtheorem{theorem}{Theorem}[section]
\newtheorem{lemma}[theorem]{Lemma}
\newtheorem{proposition}[theorem]{Proposition}
\newtheorem{corollary}[theorem]{Corollary}
\newtheorem{definition}[theorem]{Definition}
\newtheorem{remark}[theorem]{Remark}
\newcommand{\Z}{\mathbb Z}
\begin{document}
\title[Admissible bases and generic lattice ideals]{Admissible bases and generic lattice ideals}
\author{Anargyros Katsabekis}
\date{}
\maketitle
\noindent\emph{Source and licence.} Admissible bases and generic lattice ideals. Version arXiv:2607.18697v1 (CC BY 4.0). This material is distributed under the Creative Commons Attribution 4.0 International licence, http://creativecommons.org/licenses/by/4.0/, as recorded in the arXiv OAI licence metadata for this exact version and shown on the abstract page https://arxiv.org/abs/2607.18697v1. Full rights notice: licenses/arxiv-2607.18697-CC-BY-4.0.txt.\par\smallskip
\noindent\textit{Editorial scope.} Counted unit: the original \texttt{theorem} environment \texttt{thm:groebner} at source lines 882--894 together with its complete proof at lines 896--1145; the delivered document keeps source lines 359--1204 so that every referenced label and every citation resolves inside it. Native editable source is the paper's own concatenated TeX; the AMS wrapper is editorial formatting only. No publisher class, upstream script or unrelated artwork is redistributed.
\begin{lemma}\label{lem:neighbor}
Let $L\subseteq\mathbb Z^r$ be a positive lattice and let
$\mathbf u\in L\setminus\{\mathbf0\}$.
Then $\mathbf u$ is a neighbor of the origin if and only if there is
no $\mathbf b\in L\setminus\{\mathbf0,\mathbf u\}$ such that
$\mathbf b\le\mathbf u^+$.
\end{lemma}
\begin{proof}
Suppose that $\mathbf u$ is a neighbor of the origin. Then
$\{\mathbf u\}\in\Delta_L^0$, and therefore
$\{\mathbf0,\mathbf u\}\in\Delta_L$. Assume that there exists
$\mathbf b\in L\setminus\{\mathbf0,\mathbf u\}$
such that
$\mathbf b\le\mathbf u^+$.
For each $1 \leq i \leq r$, we have
$
(\mathbf0)_i=0\le(\mathbf u^+)_i,$ 
$(\mathbf u)_i\le(\mathbf u^+)_i,$ and
$(\mathbf b)_i\le(\mathbf u^+)_i.$
Since
$(\mathbf u^+)_i=\max\{0,(\mathbf u)_i\}$, either $(\mathbf0)_i$ or $(\mathbf u)_i$ is equal to
$(\mathbf u^+)_i$. Hence $
\max\{(\mathbf0)_i,(\mathbf u)_i,(\mathbf b)_i\}
=
(\mathbf u^+)_i.$
Thus
\[
\max(\{\mathbf0,\mathbf u,\mathbf b\})
=
\mathbf u^+
=
\max(\{\mathbf0,\mathbf u\}).
\]
Since
$\{\mathbf0,\mathbf u,\mathbf b\}$ and
$\{\mathbf0,\mathbf u\}$
are distinct subsets of $L$ with the same componentwise maximum, this
contradicts the fact that
$\{\mathbf0,\mathbf u\}\in\Delta_L$.
Therefore no such vector $\mathbf b$ exists.

Conversely, suppose that there is no
$\mathbf b\in L\setminus\{\mathbf0,\mathbf u\}$
satisfying
$\mathbf b\le\mathbf u^+$.
Since $L$ is positive and $\mathbf u\neq\mathbf0$, neither
$\mathbf u$ nor $-\mathbf u$ belongs to $\mathbb N^r$.
Hence
$\mathbf u\neq\mathbf u^+$.
Suppose that
$\{\mathbf0,\mathbf u\}\notin\Delta_L$.
Since
$
\max(\{\mathbf0,\mathbf u\})
=
\mathbf u^+,$
there exists a finite subset
$J\subseteq L$
such that
$J\neq\{\mathbf0,\mathbf u\}$
and
$\max(J)=\mathbf u^+$. Thus every element of $J$ is bounded coordinatewise by
$\mathbf u^+$; that is,
$\mathbf c\le\mathbf u^+$ for every $\mathbf c\in J$.
Since
$J\neq\{\mathbf0,\mathbf u\}$,
there exists
$\mathbf b\in L\setminus\{\mathbf0,\mathbf u\}$
such that
$\mathbf b\in J$.
Therefore,
$\mathbf b\le\mathbf u^+$,
contradicting the hypothesis.
Hence
$\{\mathbf0,\mathbf u\}\in\Delta_L$.
Equivalently,
$\{\mathbf u\}\in\Delta_L^0$,
so $\mathbf u$ is a neighbor of the origin.
\end{proof}
The following proposition gives the correspondence between the
neighbors of the origin and the minimal binomial generators of a
generic lattice ideal.

\begin{proposition} \label{cor:neighbors}
Let $I_L$ be a generic lattice ideal. A lattice vector
${\bf u}\in L$ is a neighbor of the origin if and only if
$
{\bf x}^{{\bf u}^{+}}-{\bf x}^{{\bf u}^{-}}$
is a minimal binomial generator of $I_L$.
\end{proposition}

\begin{proof}
By~\cite[Theorem~2.2]{Sabzrou}, for a generic lattice ideal $I_L$,
${\bf u}\in L$ is a neighbor of the origin if and only if
$
B={\bf x}^{{\bf u}^{+}}-{\bf x}^{{\bf u}^{-}}$
is an indispensable binomial of $I_L$, that is, either $B$ or $-B$
belongs to every binomial generating set of $I_L$.
By~\cite[Remark~4.4(3)]{PS}, a generic lattice ideal has a unique
minimal binomial generating set. Hence the minimal binomial generators
coincide with the indispensable binomials. Therefore
$
{\bf x}^{{\bf u}^{+}}-{\bf x}^{{\bf u}^{-}}$ is a minimal binomial generator of $I_L$ if and only if
${\bf u}$ is a neighbor of the origin.
\end{proof}
We shall also use standard facts from the theory of Gr\"obner bases.
Fix a monomial order $\prec$ on $S$. For a polynomial $f\in S$, we
denote by $\operatorname{lm}(f)$ and $\operatorname{lc}(f)$ its
leading monomial and leading coefficient, respectively, with respect
to $\prec$. Our computations rely on Buchberger's criterion together
with the coprimeness criterion.

\begin{theorem}[Buchberger's Criterion {\cite[Chapter~2, \S6, Theorem~6]{CLO}}]
\label{thm:buchberger}
A finite subset $G=\{g_1,\ldots,g_t\}\subseteq I$
is a Gr\"obner basis of an ideal $I\subseteq S$ if and only if every
$S$-polynomial $S(g_i,g_j)$ reduces to zero modulo $G$.
\end{theorem}

\begin{proposition}[{\cite[Chapter~2, Proposition~4]{CLO}}]
\label{prop:coprime}
Let $f,g\in S$. If the leading monomials of $f$ and $g$ are relatively
prime, then the $S$-polynomial $S(f,g)$ reduces to zero modulo
$\{f,g\}$.
\end{proposition}

%--------------------------------------------------------------------
\section{Admissible Lattices and the Parametric Family}
\label{sec:setup}
%--------------------------------------------------------------------
In this section we first introduce the notion of an admissible basis
for a rank-three lattice in $\mathbb Z^4$. The defining sign
conditions ensure that the lattice vectors introduced below have full
support and satisfy the coordinate properties needed in the subsequent
analysis of neighbors of the origin.

\begin{definition}\label{def:setup} A $\mathbb Z$-basis
$\{\mathbf c_1,\mathbf c_2,\mathbf c_3\}$ of a rank-three lattice
$L\subseteq\mathbb Z^4$ is called \emph{admissible} if $\mathbf c_1=(\alpha_{11},-\alpha_{12},-\alpha_{13},-\alpha_{14})$,
$\mathbf c_2=(-\alpha_{21},-\alpha_{22},-\alpha_{23},\alpha_{24})$, and
$\mathbf c_3=(-\alpha_{31},\alpha_{32},-\alpha_{33},-\alpha_{34})$,
where every $\alpha_{ij}$ is a positive integer, and the following conditions hold:
\begin{itemize}
\item[\textbf{(I)}]
$\alpha_{11}>\alpha_{21}+2\alpha_{31}$,
\item[\textbf{(II)}]
$\max(\alpha_{12},\alpha_{22})<\alpha_{32}
<\alpha_{12}+\alpha_{22}$,
\item[\textbf{(III)}]
$\alpha_{24}>\alpha_{14}+2\alpha_{34}$.
\end{itemize}

\end{definition}

The inequalities in Definition~\ref{def:setup} are chosen so that the
coordinatewise criterion of Lemma~\ref{lem:neighbor} can be applied
effectively to the vectors studied in
Section~\ref{sec:lambda}.

Throughout the remainder of the paper, given an admissible basis, we
write $\mathbf u_i=\mathbf c_i$ for $i=1,2,3$, and define
\[
\mathbf u_4=\mathbf c_1+\mathbf c_3,\;
\mathbf u_5=\mathbf c_2+\mathbf c_3,\;
\mathbf u_6=\mathbf c_1+\mathbf c_2+\mathbf c_3,\;
\mathbf u_7=\mathbf c_1+\mathbf c_2+2\mathbf c_3.
\]

The following proposition establishes a basic consequence of the
defining properties of an admissible basis that will be used
repeatedly throughout the paper.

\begin{proposition}\label{prop:full-support} Under conditions
\textup{(I)}, \textup{(II)}, and \textup{(III)}, every vector $\mathbf u_i$, $1\le i\le7$, has full support. Moreover, $(\mathbf u_i^+)_3=0$ and $(\mathbf u_i^-)_3>0$ for every $1 \leq i \leq 7$.
\end{proposition}

\begin{proof}
The third coordinates of the vectors
$\mathbf u_1,\ldots,\mathbf u_7$ are
\[
-\alpha_{13},\,
-\alpha_{23},\,
-\alpha_{33},\,
-(\alpha_{13}+\alpha_{33}),\,
-(\alpha_{23}+\alpha_{33}),\,
-(\alpha_{13}+\alpha_{23}+\alpha_{33}),\,
-(\alpha_{13}+\alpha_{23}+2\alpha_{33}).
\]
Since every $\alpha_{ij}$ is positive, each of these numbers is
negative. Hence
$
(\mathbf u_i^+)_3=0$ and $(\mathbf u_i^-)_3>0$ for every $1\le i\le7$.

It remains to verify that every $\mathbf u_i$ has full support.
The first coordinates of the vectors
$\mathbf u_1,\ldots,\mathbf u_7$ are
\[
\alpha_{11},\,
-\alpha_{21},\,
-\alpha_{31},\,
\alpha_{11}-\alpha_{31},\,
-(\alpha_{21}+\alpha_{31}),\,
\alpha_{11}-\alpha_{21}-\alpha_{31},\,
\alpha_{11}-\alpha_{21}-2\alpha_{31}.
\]
Since
$
\alpha_{11}>\alpha_{21}+2\alpha_{31}>0$
by condition~\textup{(I)}, the last expression is positive, while the
remaining quantities are clearly nonzero. Hence every first coordinate
is nonzero.

The second coordinates of the vectors
$\mathbf u_1,\ldots,\mathbf u_7$ are
\[
-\alpha_{12},\,
-\alpha_{22},\,
\alpha_{32},\,
\alpha_{32}-\alpha_{12},\,
\alpha_{32}-\alpha_{22},\,
\alpha_{32}-\alpha_{12}-\alpha_{22},\,
2\alpha_{32}-\alpha_{12}-\alpha_{22}.
\]
Condition~\textup{(II)} implies that
$\alpha_{32}-\alpha_{12}>0$, $\alpha_{32}-\alpha_{22}>0$, and 
$\alpha_{32}-\alpha_{12}-\alpha_{22}<0.$
Moreover,
\[
2\alpha_{32}-\alpha_{12}-\alpha_{22}
=
(\alpha_{32}-\alpha_{12})
+
(\alpha_{32}-\alpha_{22})
>
0.
\]
Hence each of these quantities is nonzero.

Finally, the fourth coordinates are
\[
-\alpha_{14},\,
\alpha_{24},\,
-\alpha_{34},\,
-(\alpha_{14}+\alpha_{34}),\,
\alpha_{24}-\alpha_{34},\,
\alpha_{24}-\alpha_{14}-\alpha_{34},\,
\alpha_{24}-\alpha_{14}-2\alpha_{34}.
\]
Since $ \alpha_{24}>\alpha_{14}+2\alpha_{34}$ by condition ~\textup{(III)}, the last three quantities are positive, while the first four are clearly nonzero. Hence every fourth coordinate is nonzero. Therefore each $\mathbf u_i$
has full support.
\end{proof}

Next we introduce an explicit parametric family of admissible lattices that will serve as the main source of examples throughout the remainder of the paper.

\begin{definition}\label{def:family}
For each integer $m\ge1$, let $\mathbf c_1(m)=(m+4,-(m+1),-1,-1)$,
$\mathbf c_2(m)=(-(m+1),-3,-1,m+4)$, and
$\mathbf c_3(m)=(-1,m+3,-(m+1),-1)$. Set
\[
L(m)=
\mathbb Z\mathbf c_1(m)+
\mathbb Z\mathbf c_2(m)+
\mathbb Z\mathbf c_3(m).
\]
\end{definition}


The following proposition establishes the basic properties of this
family.

\begin{proposition}\label{prop:family}
For every integer $m\ge1$, let
\[
\mathbf n(m)=
\bigl(n_1(m),n_2(m),n_3(m),n_4(m)\bigr),
\]
where
\[
\begin{aligned}
n_1(m)&=m^3+7m^2+21m+20, &
n_2(m)&=m^3+8m^2+25m+25,\\
n_3(m)&=m^3+8m^2+26m+25, &
n_4(m)&=m^3+8m^2+28m+30.
\end{aligned}
\]
Then the following hold.

\begin{enumerate}
\item[\textup{(i)}]
\[
\gcd(n_1(m),n_2(m),n_3(m),n_4(m))
=
\begin{cases}
5,&\text{if $5$ divides $m$,}\\
1,&\text{otherwise.}
\end{cases}
\]

\item[\textup{(ii)}]
The lattice $L(m)$ has admissible basis
$\{\mathbf c_1(m),\mathbf c_2(m),\mathbf c_3(m)\}$.

\item[\textup{(iii)}]
If $5$ does not divide $m$, then $
L(m)={\rm ker}_{\mathbb Z}(n_{1}(m),n_{2}(m),n_{3}(m),n_{4}(m)).$
\end{enumerate}
\end{proposition}

\begin{proof}
Let $A(m)$ be the $3\times4$ integer matrix whose rows are
$\mathbf c_1(m)$, $\mathbf c_2(m)$, and $\mathbf c_3(m)$.
For $i=1,\ldots,4$, let $\Delta_i$ denote the determinant of the
submatrix obtained by deleting the $i$-th column of $A(m)$.
A direct computation yields
\[
(-\Delta_1,\Delta_2,-\Delta_3,\Delta_4)
=
(n_1(m),n_2(m),n_3(m),n_4(m)).
\]
Hence
$
A(m)\mathbf n(m)^{\mathrm T}=\mathbf 0,$ where $\mathbf n(m)^{\mathrm T}$ denotes the transpose of
$\mathbf n(m)$.

We first prove assertion~\textup{(i)}.
Since $n_3(m)-n_2(m)=m$ and $n_4(m)-n_3(m)=2m+5$,
any common divisor of
$n_1(m),n_2(m),n_3(m),n_4(m)$
divides both $m$ and $2m+5$.
Hence
$
\gcd(n_1(m),n_2(m),n_3(m),n_4(m))$
divides $\gcd(m,2m+5)
=
\gcd(m,5).$

If $5$ does not divide $m$, then $\gcd(m,5)=1$, and therefore
\[
\gcd(n_1(m),n_2(m),n_3(m),n_4(m))=1.
\]

Assume now that $5$ divides $m$. Since every nonconstant term of $n_i(m)$ contains a factor of $m$, and the constant term is divisible by $5$, it follows that $5$ divides $n_{i}(m)$ for each $1 \leq i \leq 4$.
Hence $5$ divides
$\gcd(n_1(m),n_2(m),n_3(m),n_4(m))$.
On the other hand,
$
\gcd(n_1(m),n_2(m),n_3(m),n_4(m))$ divides
$\gcd(m,5)=5,$ so
\[
\gcd(n_1(m),n_2(m),n_3(m),n_4(m))=5.
\]

Next we prove assertion~\textup{(ii)}.
The relevant coefficients are
\[
\begin{aligned}
\alpha_{11}&=m+4,\;
\alpha_{12}=\alpha_{21}=m+1,\;
\alpha_{22}=3,\;
\alpha_{31}=1,\\
\alpha_{32}&=m+3,\;
\alpha_{14}=\alpha_{34}=1,\;
\alpha_{24}=m+4.
\end{aligned}
\]

Condition~\textup{(I)} follows from
\[
\alpha_{11}-\alpha_{21}-2\alpha_{31}
=(m+4)-(m+1)-2=1>0.
\]

For condition~\textup{(II)},
\[
\alpha_{12}+\alpha_{22}
=(m+1)+3=m+4,
\]
and therefore
\[
\alpha_{32}=m+3<m+4=\alpha_{12}+\alpha_{22}.
\]
Moreover,
\[
\max\{\alpha_{12},\alpha_{22}\}
=\max\{m+1,3\}
<m+3=\alpha_{32}.
\]

Finally,
\[
\alpha_{24}-\alpha_{14}-2\alpha_{34}
=(m+4)-1-2=m+1>0,
\]
so condition~\textup{(III)} is satisfied.

Finally, we prove assertion~\textup{(iii)}.
Since $A(m)$ has rank~$3$, its Smith normal form is
${\rm diag}(e_1,e_2,e_3,0),$ where $e_1,e_2,e_3$ are the invariant factors.
The discussion following Theorem~II.9 in
\cite{Newman} shows that
$
e_1e_2e_3
=
\gcd(n_1(m),n_2(m),n_3(m),n_4(m)).$

If $5$ does not divide $m$, then assertion~\textup{(i)} gives
$e_1e_2e_3=1$.
Since $e_1$ divides $e_2$ and $e_2$ divides $e_3$, we obtain
$e_1=e_2=e_3=1$.
By the fundamental theorem of finitely generated abelian groups
(\cite[Section~12.1]{DF}),
\[
\mathbb Z^4/L(m)
\cong
\mathbb Z
\oplus
\mathbb Z/e_1\mathbb Z
\oplus
\mathbb Z/e_2\mathbb Z
\oplus
\mathbb Z/e_3\mathbb Z,
\]
and therefore $\mathbb Z^4/L(m)$ is torsion free.
Moreover,
$
A(m)\mathbf n(m)^{\mathrm T}={\bf 0},$
so
$$
L(m) \subseteq {\rm ker}_{\mathbb Z}(n_{1}(m),n_{2}(m),n_{3}(m),n_{4}(m)).$$
Suppose that $L(m) \neq {\rm ker}_{\mathbb Z}(n_{1}(m),n_{2}(m),n_{3}(m),n_{4}(m))$. Since both lattices have
rank~$3$, the quotient
${\rm ker}_{\mathbb Z}(n_{1}(m),n_{2}(m),n_{3}(m),n_{4}(m))/L(m)$
is a nontrivial finite subgroup of $\mathbb Z^4/L(m)$. Since every element of a finite group has finite order, it follows that $\mathbb Z^4/L(m)$ contains a nontrivial torsion element, contradicting the fact that
$\mathbb Z^4/L(m)$ is torsion-free. Hence
$
L(m)={\rm ker}_{\mathbb Z}(n_{1}(m),n_{2}(m),n_{3}(m),n_{4}(m)).$
\end{proof}

The next proposition establishes that the lattices $L(m)$ are
positive.

\begin{proposition}
\label{prop:positive}
For every integer $m\ge1$, the lattice $L(m)$ is positive.
\end{proposition}

\begin{proof}
Let $\mathbf u=(u_1,\ldots,u_4)\in L(m)\cap\mathbb N^4$. Since
$A(m)\mathbf n(m)^{\mathrm T}={\bf 0}$, the vector $\mathbf n(m)$ is
orthogonal to every vector of $L(m)$. Hence
$
\mathbf n(m)\cdot\mathbf u=0.$
Since each coordinate of $\mathbf n(m)$ is positive and each
$u_i\ge0$, we have
\[
\mathbf n(m)\cdot\mathbf u
=
n_1(m)u_1+\cdots+n_4(m)u_4,
\]
where every summand is nonnegative. Therefore
$\mathbf n(m)\cdot\mathbf u=0$ if and only if
$u_1=u_2=u_3=u_4=0$. Hence
$\mathbf u=\mathbf0$, and consequently
$
L(m)\cap\mathbb N^4=\{\mathbf0\}.$
Thus $L(m)$ is positive.
\end{proof}

\begin{remark}
For $m=1$, Proposition~\ref{prop:family} gives
$\mathbf n(1)=(49,59,60,67).$ The affine monomial curve parametrized by
$x_1=t^{49}$, $x_2=t^{59}$, $x_3=t^{60}$, and $x_4=t^{67}$
is discussed in Example~4.3 of~\cite{OP},
where they observe that the lattice vectors
corresponding to the minimal binomial generators of $I_{L(1)}$
cannot be arranged as the vertices of the three-dimensional unit cube.
\end{remark}

Throughout this section we use the graded reverse lexicographic order
$\prec$ determined by the grading
$
\deg(x_i)=n_i(m)$, for $1\le i\le4,$
and the variable order
$
x_1\succ x_2\succ x_4\succ x_3.$
Thus, for exponent vectors
$\mathbf a=(a_1,a_2,a_3,a_4)$ and
$\mathbf b=(b_1,b_2,b_3,b_4)$,
we have
\[
{\bf x}^{\mathbf a}\prec {\bf x}^{\mathbf b}
\quad\text{if}\quad
\sum_{i=1}^4 a_i n_i(m)
<
\sum_{i=1}^4 b_i n_i(m),
\]
or if
\[
\sum_{i=1}^4 a_i n_i(m)
=
\sum_{i=1}^4 b_i n_i(m),
\]
and the last nonzero coordinate of
$
(a_1-b_1,\;a_2-b_2,\;a_4-b_4,\;a_3-b_3)$
is positive.


The following notation will be used throughout the remainder of this
section. Let
\[
\begin{aligned}
g_1&=x_1^{m+4}-x_2^{m+1}x_3x_4, \;
g_2=x_4^{m+4}-x_1^{m+1}x_2^3x_3, \;
g_3=x_2^{m+3}-x_1x_3^{m+1}x_4,\\
g_4&=x_1^{m+3}x_2^2-x_3^{m+2}x_4^2, \;
g_5=x_2^{m}x_4^{m+3}-x_1^{m+2}x_3^{m+2}, \;
g_6=x_1^2x_4^{m+2}-x_2x_3^{m+3},\\
g_7&=x_1x_2^{m+2}x_4^{m+1}-x_3^{2m+4}.
\end{aligned}
\]



\begin{theorem}\label{thm:groebner}
For every $m\ge1$, the set
$
G=\{g_1,\ldots,g_7\}$
is the reduced Gr\"obner basis of the lattice ideal $I_{L(m)}$ with
respect to $\prec$. Moreover, $I_{L(m)}$ is generic. The set $G$ is the unique minimal
system of binomial generators of $I_{L(m)}$, and
\[
N(L(m))
=
\{\pm\mathbf u_i\mid1\le i\le7\}.
\]
\end{theorem}

\begin{proof} We first show that $G$ is a Gr\"obner basis for the ideal $J=\langle g_{i} \mid 1 \leq i \leq 7 \rangle$ with
respect to $\prec$. Since every element of $G$ is a binomial
$\mathbf x^{\mathbf u^+}-\mathbf x^{\mathbf u^-}$ with
$\mathbf u=\mathbf u^+-\mathbf u^-\in L(m)$, we have $J\subseteq I_{L(m)}$. The leading monomials of the elements of $G$ are
\[
\operatorname{lm}(g_1)=x_1^{m+4},\;
\operatorname{lm}(g_2)=x_4^{m+4},\;
\operatorname{lm}(g_3)=x_2^{m+3},\;
\operatorname{lm}(g_4)=x_1^{m+3}x_2^2, \]
\[\operatorname{lm}(g_5)=x_2^mx_4^{m+3},\;
\operatorname{lm}(g_6)=x_1^2x_4^{m+2},\;
\operatorname{lm}(g_7)=x_1x_2^{m+2}x_4^{m+1}.
\]

We verify Buchberger's criterion. Since the leading monomials of $g_i$ and $g_j$ are relatively prime for
each of the pairs
\[
\{g_1,g_2\},\;
\{g_1,g_3\},\;
\{g_1,g_5\},\;
\{g_2,g_3\},\;
\{g_2,g_4\},\;
\{g_3,g_6\},
\]
their $S$-polynomials reduce to zero by
Proposition~\ref{prop:coprime}. For the following pairs, we have
\[
\begin{aligned}
S(g_1,g_4)&=x_2^2g_1-x_1g_4=-x_3x_4g_3\stackrel{g_3}{\longrightarrow}0,\\
S(g_1,g_6)&=x_4^{m+2}g_1-x_1^{m+2}g_6=-x_2x_3g_5\stackrel{g_5}{\longrightarrow}0,\\
S(g_2,g_5)&=x_2^mg_2-x_4g_5=-x_1^{m+1}x_3g_3\stackrel{g_3}{\longrightarrow}0,\\
S(g_2,g_6)&=x_1^2g_2-x_4^2g_6=-x_2x_3g_4\stackrel{g_4}{\longrightarrow}0,\\
S(g_3,g_4)&=x_1^{m+3}g_3-x_2^{m+1}g_4=-x_3^{m+1}x_4g_1\stackrel{g_1}{\longrightarrow}0,\\
S(g_3,g_5)&=x_4^{m+3}g_3-x_2^3g_5=-x_1x_3^{m+1}g_2\stackrel{g_2}{\longrightarrow}0,\\
S(g_3,g_7)&=x_1x_4^{m+1}g_3-x_2g_7=-x_3^{m+1}g_6\stackrel{g_6}{\longrightarrow}0,\\
S(g_4,g_6)&=x_4^{m+2}g_4-x_1^{m+1}x_2^2g_6=-x_3^{m+2}g_2\stackrel{g_2}{\longrightarrow}0,\\
S(g_4,g_7)&=x_2^mx_4^{m+1}g_4-x_1^{m+2}g_7=-x_3^{m+2}g_5\stackrel{g_5}{\longrightarrow}0,\\
S(g_5,g_6)&=x_1^2g_5-x_2^mx_4g_6=-x_3^{m+2}g_1\stackrel{g_1}{\longrightarrow}0,\\
S(g_5,g_7)&=x_1x_2^2g_5-x_4^2g_7=-x_3^{m+2}g_4\stackrel{g_4}{\longrightarrow}0,\\
S(g_6,g_7)&=x_2^{m+2}g_6-x_1x_4g_7=-x_3^{m+3}g_3\stackrel{g_3}{\longrightarrow}0.
\end{aligned}
\]

It remains to consider the pairs
$\{g_1,g_7\}$,
$\{g_2,g_7\}$, and
$\{g_4,g_5\}$.

\begin{enumerate}
\item[\textup{(i)}] \emph{The pair $\{g_1,g_7\}$.} Since
$
\operatorname{lcm}\!\left(\operatorname{lm}(g_1),\operatorname{lm}(g_7)\right)
=
x_1^{m+4}x_2^{m+2}x_4^{m+1},$
we have
\[
S(g_1,g_7)
=
x_2^{m+2}x_4^{m+1}g_1
-
x_1^{m+3}g_7.
\]
Hence
\[
\begin{aligned}
S(g_1,g_7)
&=
x_2^{m+2}x_4^{m+1}
(x_1^{m+4}-x_2^{m+1}x_3x_4)
-
x_1^{m+3}
(x_1x_2^{m+2}x_4^{m+1}-x_3^{2m+4})\\
&=
x_1^{m+3}x_3^{2m+4}
-
x_2^{2m+3}x_3x_4^{m+2}.
\end{aligned}
\]
The leading monomial of $S(g_1,g_7)$ is
$x_2^{2m+3}x_3x_4^{m+2}$, which is divisible by
$\operatorname{lm}(g_3)$. We have
\[S(g_1,g_7)
\stackrel{g_3}{\longrightarrow}
x_1^{m+3}x_3^{2m+4}-x_1x_2^mx_3^{m+2}x_4^{m+3}=
-x_1x_3^{m+2}g_5
\stackrel{g_5}{\longrightarrow}
0.\]

\item[\textup{(ii)}] \emph{The pair $\{g_2,g_7\}$.} Since
$
\operatorname{lcm}\!\left(\operatorname{lm}(g_2),\operatorname{lm}(g_7)\right)
=
x_1x_2^{m+2}x_4^{m+4},$
we have
\[
S(g_2,g_7)
=
x_1x_2^{m+2}g_2
-
x_4^3g_7.
\]
Hence
\[
S(g_2,g_7)
=
x_3^{2m+4}x_4^3
-
x_1^{m+2}x_2^{m+5}x_3.
\]
The leading monomial of $S(g_2,g_7)$ is
$x_1^{m+2}x_2^{m+5}x_3$, which is divisible by
$\operatorname{lm}(g_3)$. We have
\[S(g_2,g_7)
\stackrel{g_3}{\longrightarrow}
x_3^{2m+4}x_4^3-x_1^{m+3}x_2^2x_3^{m+2}x_4=
-x_3^{m+2}x_4\,g_4
\stackrel{g_4}{\longrightarrow}
0.\]

\item[\textup{(iii)}] \emph{The pair $\{g_4,g_5\}$.} We distinguish two cases.

\begin{enumerate}
\item[\textup{(a)}] \emph{$m\ge2$.} Since
$
\operatorname{lcm}\!\left(\operatorname{lm}(g_4),\operatorname{lm}(g_5)\right)
=
x_1^{m+3}x_2^mx_4^{m+3},$
we have
\[
S(g_4,g_5)
=
x_2^{m-2}x_4^{m+3}g_4
-
x_1^{m+3}g_5.
\]
Hence
\[
S(g_4,g_5)
=
x_1^{2m+5}x_3^{m+2}
-
x_2^{m-2}x_3^{m+2}x_4^{m+5}.
\]
The leading monomial of $S(g_4,g_5)$ is
$x_1^{2m+5}x_3^{m+2}$, which is divisible by
$\operatorname{lm}(g_1)$. We have
$$S(g_4,g_5)
\stackrel{g_1}{\longrightarrow}
x_1^{m+1}x_2^{m+1}x_3^{m+3}x_4
-
x_2^{m-2}x_3^{m+2}x_4^{m+5}=
-x_2^{m-2}x_3^{m+2}x_4\,g_2
\stackrel{g_2}{\longrightarrow}
0.$$

\item[\textup{(b)}] \emph{$m=1$.} Since
$
\operatorname{lcm}\!\left(\operatorname{lm}(g_4),\operatorname{lm}(g_5)\right)
=
x_1^4x_2^2x_4^4,$
we have
\[
S(g_4,g_5)
=
x_4^4g_4
-
x_1^4x_2g_5.
\]
Hence
\[
S(g_4,g_5)
=
x_1^7x_2x_3^3
-
x_3^3x_4^6.
\]
The leading monomial of $S(g_4,g_5)$ is
$x_1^7x_2x_3^3$, which is divisible by
$\operatorname{lm}(g_1)$. We have
\[S(g_4,g_5)
\stackrel{g_1}{\longrightarrow}
x_1^2x_2^3x_3^4x_4
-
x_3^3x_4^6=
-x_3^3x_4\,g_2
\stackrel{g_2}{\longrightarrow}
0.
\]

\end{enumerate}

\end{enumerate}
Therefore every $S$-polynomial reduces to zero modulo $G$. By
Buchberger's criterion (Theorem~\ref{thm:buchberger}), $G$ is a
Gr\"obner basis of $J$.

Next we show that $G$ is the reduced Gr\"obner basis of $J$ with respect to $\prec$. Since $\operatorname{lc}(g_i)=1$ for every $1\le i\le7$, it remains to
verify that no monomial of an element of $G$ is divisible by the
leading monomial of another element of $G$. This is immediate, since
no leading monomial divides another leading monomial, and every
non-leading monomial contains a positive power of $x_3$, whereas no
leading monomial does. Hence $G$ is the reduced Gr\"obner basis of
$J$. 

Since $x_3$ does not divide any element of $G$, the set $G$ is a Gr\"obner basis of
$J:x_3^{\infty}$ by \cite[Lemma 12.1]{Stu}, thus $J:x_3^{\infty}=J$.

We now show that $J=I_{L(m)}$. It is well known that there is a positive integer $r$ such that $$J=(J:(x_1x_2x_4)^\infty) \cap (J+(x_1x_2x_4)^r).$$ We have \[
x_3^{r(2m+4)}
=
x_2^{(m+1)r}x_4^{mr}\,(x_1x_2x_4)^r
\;-\;
g_7\sum_{i=0}^{r-1}x_3^{(2m+4)(r-1-i)}\bigl(x_1x_2^{m+2}x_4^{m+1}\bigr)^i.
\] Thus $x_3^{r(2m+4)}\in J+(x_1x_2x_4)^r$, and therefore
$\bigl(J+(x_1x_2x_4)^r\bigr):x_3^\infty=S$. Hence
\[
J
=
J:x_3^\infty
=
\Bigl[\bigl(J:(x_1x_2x_4)^\infty\bigr)
\cap
\bigl(J+(x_1x_2x_4)^r\bigr)\Bigr]:x_3^\infty,
\] so \[
J
=
\bigl(J:(x_1x_2x_3x_4)^\infty\bigr)\cap \bigl(J+(x_1x_2x_4)^r\bigr):x_3^\infty
=\bigl(J:(x_1x_2x_3x_4)^\infty\bigr)\cap S=
J:(x_1x_2x_3x_4)^\infty.
\] By \cite[Corollary 2.5]{ES}, we have that $J=I_{L(m)}$. Since each binomial $g_i=\mathbf{x}^{\mathbf{u}_i^+}-\mathbf{x}^{\mathbf{u}_i^-}$
has full support, it follows that $I_{L(m)}$ is generic.

Moreover, no monomial occurring in an element of $G$ divides another
such monomial, and each monomial occurs in exactly one element of $G$.
Hence every monomial occurring in $G$ is indispensable, that is, every
system of binomial generators of $I_{L(m)}$ contains a binomial having
this monomial as one of its terms
\cite[Remark~2.3]{CTV}. Furthermore, the graph associated with $G$ in
the sense of Theorem~3.3 of~\cite{CTV} consists of seven connected
components, each of which is a single edge joining the two monomials
of one of the binomials $g_1,\ldots,g_7$. Therefore every binomial
$g_i$ is indispensable by \cite[Theorem~3.3]{CTV}. Since
$I_{L(m)}$ is generic, \cite[Remark~4.4(3)]{PS}
implies that $G=\{g_1,\ldots,g_7\}$ is the unique minimal system of
binomial generators of $I_{L(m)}$. 

Finally, Proposition~\ref{cor:neighbors} yields
$
N(L(m))=\{\pm\mathbf u_1,\ldots,\pm\mathbf u_7\}.$
\end{proof}


We now determine the $f$-vector of the algebraic Scarf complex of
$I_{L(m)}$. Since $I_{L(m)}$ is generic, \cite[Theorem~4.2]{PS} shows that the algebraic Scarf complex of
$I_{L(m)}$ is the minimal free resolution of $S/I_{L(m)}$.


\begin{proposition}\label{prop:scarf}
For every integer $m\ge1$, the ring $S/I_{L(m)}$ has Betti numbers
$$
\beta_0(S/I_{L(m)})=1, \ \beta_1(S/I_{L(m)})=7, \ \beta_2(S/I_{L(m)})=12, \ \beta_3(S/I_{L(m)})=6,$$ and
$
\beta_i(S/I_{L(m)})=0$ for every $i \geq 4$. Consequently, the algebraic Scarf complex of $I_{L(m)}$ has
$f$-vector
$
(f_0,f_1,f_2)=(7,12,6).$
\end{proposition}
 
\begin{proof} By Proposition~\ref{prop:full-support} and
Theorem~\ref{thm:groebner}, the lattice ideal $I_{L(m)}$ is generic
and has exactly seven minimal binomial generators.

The proof of Corollary~1.3 in~\cite{Ojeda} shows that, for a generic
lattice ideal $I_L$ of codimension~$3$ with $t\ge7$ minimal binomial
generators,
\[
\beta_0(S/I_L)=1, \
\beta_1(S/I_L)=t, \
\beta_2(S/I_L)=2t-2, \
\beta_3(S/I_L)=t-1,
\]
and
$
\beta_i(S/I_L)=0$ for every $i \geq 4$.
Since $I_{L(m)}$ has $t=7$ minimal generators, it follows that
\[
\beta_0(S/I_{L(m)})=1, \
\beta_1(S/I_{L(m)})=7, \
\beta_2(S/I_{L(m)})=12, \
\beta_3(S/I_{L(m)})=6,
\]
and
$
\beta_i(S/I_{L(m)})=0$ for every $i \geq 4$.

By~\cite[Theorem~4.2]{PS}, the minimal free resolution of
$S/I_{L(m)}$ is supported on the algebraic Scarf complex.
Hence the numbers of vertices, edges, and triangles of the algebraic
Scarf complex are
$\beta_1(S/I_{L(m)})$,
$\beta_2(S/I_{L(m)})$, and
$\beta_3(S/I_{L(m)})$, respectively.
Therefore its $f$-vector is $
(f_0,f_1,f_2)=(7,12,6).$
\end{proof}


\section{Restrictions on neighbors of the origin}
\label{sec:lambda}


\begin{thebibliography}{99}
\bibitem[CLO]{CLO} D.~Cox, J.~Little, and D.~O'Shea, \textit{Ideals, Varieties, and Algorithms: An Introduction to Computational Algebraic Geometry and Commutative Algebra}, Undergraduate Texts in Mathematics, Springer, Cham, 2015.
\bibitem[CTV]{CTV} H.~Charalambous, A.~Thoma, and M.~Vladoiu, \textit{Binomial fibers and indispensable binomials}, J.\ Symb.\ Comput.\ \textbf{74} (2016), 578--591.
\bibitem[DF]{DF} D.~Dummit and R.~Foote, \textit{Abstract Algebra}, 3rd ed., John Wiley \& Sons, Hoboken, NJ, 2004.
\bibitem[ES]{ES} D.~Eisenbud and B.~Sturmfels, \textit{Binomial ideals}, Duke Math.\ J.\ \textbf{84} (1996), no.~1, 1--45.
\bibitem[Newman]{Newman} M.~Newman, \textit{Integral Matrices}, Pure and Applied Mathematics, vol.~45, Academic Press, New York, 1972.
\bibitem[OP]{OP} I.~Ojeda and P.~Pis\'on-Casares, \textit{On the hull resolution of an affine monomial curve}, J.~Pure Appl.\ Algebra \textbf{192} (2004), 53--67.
\bibitem[Ojeda]{Ojeda} I.~Ojeda, \textit{Examples of generic lattice ideals of codimension~3}, Comm.\ Algebra \textbf{36} (2008), no.~1, 279--287.
\bibitem[PS]{PS} I.~Peeva and B.~Sturmfels, \textit{Generic lattice ideals}, J.~Amer.\ Math.\ Soc.\ \textbf{11} (1998), no.~2, 363--373.
\bibitem[Sabzrou]{Sabzrou} H.~Sabzrou, \textit{Scarf lattice ideals}, Portugaliae Math.\ \textbf{68} (2011), no.~4, 369--380.
\bibitem[Stu]{Stu} B.~Sturmfels, \textit{Gr\"obner Bases and Convex Polytopes}, University Lecture Series, vol.~8, Amer.\ Math.\ Soc., Providence, RI, 1996.
\end{thebibliography}
\end{document}
